\begin{ex}
Số La Mã V biểu thị số tự nhiên nào?
\shortans{5}
\loigiai{
V là 5.
}
\end{ex}

\begin{ex}
Số La Mã X biểu thị số tự nhiên nào?
\shortans{10}
\loigiai{
X là 10.
}
\end{ex}

\begin{ex}
Số La Mã I biểu thị số tự nhiên nào?
\shortans{1}
\loigiai{
I là 1.
}
\end{ex}

\begin{ex}
Số La Mã IV biểu thị số tự nhiên nào?
\shortans{4}
\loigiai{
IV là 4.
}
\end{ex}

\begin{ex}
Số La Mã VI biểu thị số tự nhiên nào?
\shortans{6}
\loigiai{
VI là 6.
}
\end{ex}

\begin{ex}
Số La Mã IX biểu thị số tự nhiên nào?
\shortans{9}
\loigiai{
IX là 9.
}
\end{ex}

\begin{ex}
Số La Mã XI biểu thị số tự nhiên nào?
\shortans{11}
\loigiai{
XI là 11.
}
\end{ex}

\begin{ex}
Số La Mã XII biểu thị số tự nhiên nào?
\shortans{12}
\loigiai{
XII là 12.
}
\end{ex}

\begin{ex}
Số La Mã XIV biểu thị số tự nhiên nào?
\shortans{14}
\loigiai{
XIV là 14.
}
\end{ex}

\begin{ex}
Số La Mã XV biểu thị số tự nhiên nào?
\shortans{15}
\loigiai{
XV là 15.
}
\end{ex}

\begin{ex}
Số La Mã XIX biểu thị số tự nhiên nào?
\shortans{19}
\loigiai{
XIX là 19.
}
\end{ex}

\begin{ex}
Số La Mã XX biểu thị số tự nhiên nào?
\shortans{20}
\loigiai{
XX là 20.
}
\end{ex}

\begin{ex}
Số La Mã XXI biểu thị số tự nhiên nào?
\shortans{21}
\loigiai{
XXI là 21.
}
\end{ex}

\begin{ex}
Số La Mã VII biểu thị số tự nhiên nào?
\shortans{7}
\loigiai{
VII là 7.
}
\end{ex}

\begin{ex}
Số La Mã VIII biểu thị số tự nhiên nào?
\shortans{8}
\loigiai{
VIII là 8.
}
\end{ex}

\begin{ex}
Số La Mã XIII biểu thị số tự nhiên nào?
\shortans{13}
\loigiai{
XIII là 13.
}
\end{ex}

\begin{ex}
Số La Mã XVI biểu thị số tự nhiên nào?
\shortans{16}
\loigiai{
XVI là 16.
}
\end{ex}

\begin{ex}
Số La Mã XVII biểu thị số tự nhiên nào?
\shortans{17}
\loigiai{
XVII là 17.
}
\end{ex}

\begin{ex}
Số La Mã XVIII biểu thị số tự nhiên nào?
\shortans{18}
\loigiai{
XVIII là 18.
}
\end{ex}

\begin{ex}
Đồng hồ La Mã có kim ngắn chỉ vào số IX, kim dài chỉ vào số XII. Hỏi đồng hồ chỉ mấy giờ đúng?
\shortans{9}
\loigiai{
IX là 9, đồng hồ chỉ 9 giờ đúng.
}
\end{ex}

\begin{ex}
Đồng hồ La Mã có kim ngắn chỉ vào số IV, kim dài chỉ vào số XII. Đồng hồ chỉ mấy giờ đúng?
\shortans{4}
\loigiai{
IV là 4, đồng hồ chỉ 4 giờ đúng.
}
\end{ex}

\begin{ex}
Tính giá trị bằng số tự nhiên: V + VII = \dots
\shortans{12}
\loigiai{
$5 + 7 = 12$.
}
\end{ex}

\begin{ex}
Tính giá trị bằng số tự nhiên: X + IX = \dots
\shortans{19}
\loigiai{
$10 + 9 = 19$.
}
\end{ex}

\begin{ex}
Tính giá trị bằng số tự nhiên: XX - VIII = \dots
\shortans{12}
\loigiai{
$20 - 8 = 12$.
}
\end{ex}

\begin{ex}
Dùng que tính để xếp số La Mã VII cần bao nhiêu que tính?
\shortans{4}
\loigiai{
Chữ V cần 2 que, mỗi chữ I cần 1 que: $2 + 1 + 1 = 4$ que tính.
}
\end{ex}
